\documentclass[conference]{IEEEtran}
\usepackage{cite}
\usepackage{amsmath,amssymb,amsfonts}
\usepackage{algorithmic}
\usepackage{graphicx}
\usepackage{textcomp}
\usepackage{xcolor}
\usepackage{booktabs}
\usepackage{url}
\usepackage{hyperref}

\begin{document}

\title{recovery\_eval: State-Matched and Provenance-Aware Evaluation of Recovery Behavior in Language-Model Reasoning}

\author{\IEEEauthorblockN{Sham Satish Thakare}
\IEEEauthorblockA{\textit{Independent Researcher} \\
Pune, Maharashtra, India \\
Email: shamthakare3000@gmail.com}
}

\maketitle

\begin{abstract}
Evaluating whether different language-model policy configurations exhibit improved recovery from intermediate reasoning errors is difficult because overall continuation quality can confound recovery-specific behavior. We introduce recovery\_eval, a reproducible data-centric evaluation framework that compares verifier-defined recovery states with prospectively matched reference-control states using frozen structural covariates. The framework combines state matching with append-only exposure governance, primitive neural-rollout provenance, and independent analysis reconstruction. We apply recovery\_eval to 400 genuine neural continuations generated by two released Qwen2.5-Math 1.5B checkpoint-interface configurations on 20 prospectively isolated GSM8K evaluation problems. The Instruct checkpoint showed higher observed continuation success than the Base checkpoint in both recovery and matched-control states, but the improvement was larger for controls. The resulting matched recovery-specific contrast was minus 0.110, with a descriptive 95 percent problem-level bootstrap interval from minus 0.240 to 0.030. Thus, under this evaluation, we did not observe evidence of a recovery-specific advantage for the Instruct checkpoint. The result illustrates how aggregate benchmark improvements can obscure state-specific reasoning behavior and motivates provenance-aware, state-matched evaluation for language-model diagnostics.
\end{abstract}

\begin{IEEEkeywords}
Large Language Models, Mathematical Reasoning, Model Evaluation, Data-Centric AI, Reproducibility
\end{IEEEkeywords}

\section{Introduction}
Large language models (LLMs) trained on multi-step mathematical reasoning demonstrate remarkable capacity for generating complex chain-of-thought solutions \cite{cobbe2021gsm8k, qwen25math2024}. A foundational question in AI evaluation is whether different released language-model configurations exhibit recovery behavior beyond their general continuation-quality differences \cite{lightman2023process, zelikman2022star}.

Evaluating error-recovery capability is fraught with statistical confounding. Standard benchmark evaluation measures aggregate end-to-end accuracy, which conflates baseline generation fluency with true error recovery. Simply comparing model rollouts from error-containing prefixes against valid prefixes ignores critical confounding covariates, including trajectory depth, remaining solution complexity, problem difficulty, and token length. A post-trained model may exhibit higher accuracy on error prefixes simply because its overall baseline generation quality has improved across all prompt conditions, rather than possessing a specialized error-recovery mechanism.

To address these diagnostic challenges, we introduce \texttt{recovery\_eval}, a data-centric evaluation infrastructure for measuring state-specific recovery behavior in reasoning LLMs. The primary contributions of this work are threefold:
\begin{enumerate}
    \item \textbf{Verifier-Defined State Matching Protocol}: A prospective matching methodology that pairs verifier-identified error states with valid reference-control states using frozen structural covariates and an exact matching constraint on categorical problem attributes.
    \item \textbf{Provenance and Exposure Governance Architecture}: An append-only evidence ledger that locks evaluation items, logs primitive BPE token-level generation continuations, verifies model weight manifests, and enables 100\% independent analytical reconstruction.
    \item \textbf{Genuine Two-Checkpoint Framework Demonstration}: A 400-rollout empirical evaluation across \texttt{Qwen2.5-Math-1.5B} Base and Instruct checkpoints, demonstrating that aggregate continuation success gains (+0.4300 on recovery vs +0.5400 on control) do not automatically translate into a detectable recovery-specific advantage ($D_{\text{recovery}} = -0.1100$).
\end{enumerate}

\section{Related Work}
Recent research has focused on reasoning self-correction, test-time compute scaling, and process reward modeling \cite{snell2024scaling, wang2022selfconsistency}. Table \ref{tab:related_work} compares \texttt{recovery\_eval} with existing evaluation methodologies.

\begin{table}[htbp]
\caption{Comparison with Existing Reasoning Evaluation Frameworks}
\label{tab:related_work}
\centering
\begin{tabular}{p{2.2cm}p{1.4cm}p{1.6cm}p{1.8cm}}
\toprule
\textbf{Framework / Work} & \textbf{State Matching} & \textbf{Exposure Ledger} & \textbf{Primitive Provenance} \\
\midrule
End-to-End GSM8K \cite{cobbe2021gsm8k} & None & No & Final Answer Only \\
PRM800K \cite{lightman2023process} & Unmatched & Partial & Step Scores \\
STaR \cite{zelikman2022star} & Unmatched & No & Filtered Traces \\
\textbf{\texttt{recovery\_eval} (Ours)} & \textbf{Prospectively Matched} & \textbf{Append-Only LEDGER} & \textbf{Full BPE Token + Raw JSONL} \\
\bottomrule
\end{tabular}
\end{table}

Existing benchmark designs generally do not explicitly control for the structural state differences targeted by our matched recovery/control protocol.

\section{Problem Formulation}
Let $s$ represent a partial reasoning trajectory prefix for a mathematical word problem $p \in \mathcal{P}$. A deterministic verifier $\mathcal{V}(s) \in \{0, 1\}$ evaluates whether prefix $s$ contains any logical or arithmetic errors.

A \textit{recovery state} $s_R$ is a partial prefix containing a verifier-identified intermediate error. A \textit{reference-control state} $s_C$ is a valid intermediate prefix for the same problem instance at an equivalent stage of resolution.

For a target policy $\pi$ and baseline policy $\pi_0$, let $V(s) \in \{0, 1\}$ denote the binary continuation success indicator of completing the solution correctly from prefix $s$. We define the matched recovery-specific contrast $D_{\text{recovery}}$ as:
\begin{equation}
D_{\text{recovery}} = \mathbb{E}[V_{\pi}(s_R) - V_{\pi_0}(s_R)] - \mathbb{E}[V_{\pi}(s_C) - V_{\pi_0}(s_C)]
\end{equation}

\section{The \texttt{recovery\_eval} Framework}
The \texttt{recovery\_eval} framework provides an end-to-end evaluation pipeline comprising three main modules: state generation, prospective matching, and governance-backed neural rollout execution.

\subsection{System Architecture}
The framework operates as a decoupled pipeline: 1) Registry Phase: Evaluation items and candidate prefixes are locked in an append-only registry; 2) Matching Phase: Recovery states are prospectively matched to reference control states prior to inference; 3) Execution Phase: Continuations are sampled under sealed generation settings and evaluated by the primitive verifier.

\section{State Construction and Prospective Matching}
Recovery states $s_R$ are constructed by introducing controlled single-step arithmetic perturbations into valid intermediate reasoning prefixes. Reference control states $s_C$ maintain valid prefix steps.

To prevent confounding due to trajectory length or problem complexity, each recovery state $s_R$ is matched to a control state $s_C$ using a normalized weighted-L1 Manhattan distance over continuous structural covariates:
\begin{equation}
d(i, j) = \sum_{k=1}^K w_k \frac{|x_{ik} - x_{jk}|}{s_k}
\end{equation}
subject to an exact categorical matching constraint on reasoning operation type and problem difficulty. Table \ref{tab:covariates} lists the prospective matching parameters.

\begin{table}[htbp]
\caption{Frozen Prospective Matching Covariates and Scales}
\label{tab:covariates}
\centering
\begin{tabular}{lccc}
\toprule
\textbf{Covariate Name} & \textbf{Type} & \textbf{Weight ($w_k$)} & \textbf{Scale ($s_k$)} \\
\midrule
Trajectory Depth & Continuous & 0.4 & 1.5 \\
Remaining Solution Length & Continuous & 0.4 & 1.0 \\
Token Length & Continuous & 0.2 & 15.0 \\
Reasoning Operation & Categorical & Exact Match & -- \\
Problem Difficulty & Categorical & Exact Match & -- \\
\bottomrule
\end{tabular}
\end{table}

\section{Provenance and Exposure Governance}
To guarantee artifact integrity and prevent item re-selection or data leakage, \texttt{recovery\_eval} implements an append-only event ledger. Every item exposure event is logged with cryptographic parent hashing \cite{rosenbaum1983central, ho2007matching}.

Each raw rollout record in \texttt{RAW\_NEURAL\_ROLLOUTS.jsonl} captures input token IDs, generated BPE token continuation IDs, monotonic generation timestamps, resolved HuggingFace commit hashes, weight manifest SHA-256 digests, and deterministic verifier outputs.

\section{Experimental Design}
We evaluate the framework across 20 fresh, prospectively isolated GSM8K test items ($N=20$). For each problem item, 1 recovery state and 1 matched control state are evaluated across 2 policy configurations (\texttt{Qwen2.5-Math-1.5B} Base and Instruct) using 5 stochastic rollout seeds ($S=5$), yielding exactly 400 continuations ($20 \times 2 \times 2 \times 5 = 400$).

Table \ref{tab:provenance} details the execution design and provenance metrics.

\begin{table}[htbp]
\caption{Evaluation Design and Provenance Statistics}
\label{tab:provenance}
\centering
\begin{tabular}{lp{5.2cm}}
\toprule
\textbf{Attribute} & \textbf{Specification / Value} \\
\midrule
Base Model & \texttt{Qwen/Qwen2.5-Math-1.5B} (\texttt{4a83ca6e}) \\
Instruct Model & \texttt{Qwen/Qwen2.5-Math-1.5B-Instruct} (\texttt{aafeb0fc}) \\
Hardware Device & Apple Silicon MPS (\texttt{mps:0}) \\
Total Rollouts & 400 (200 Base, 200 Instruct) \\
Total Generated Tokens & 19,212 BPE tokens \\
Measured Duration & 1,755.86 seconds ($\approx 29.26$ minutes) \\
Base Throughput & 11.86 tokens/sec (11,354 tokens in 957.37s) \\
Instruct Throughput & 9.84 tokens/sec (7,858 tokens in 798.49s) \\
BPE Decode Match & 100.0\% (400/400 exact match) \\
Raw JSONL SHA-256 & \texttt{51b5a157d9e44102caeb86d0b356f558...} \\
\bottomrule
\end{tabular}
\end{table}

\section{Results}
Table \ref{tab:outcomes} presents continuation success rates across states and policy configurations.

\begin{table}[htbp]
\caption{Observed Empirical Continuation Outcomes and Contrast}
\label{tab:outcomes}
\centering
\begin{tabular}{lccc}
\toprule
\textbf{State Condition} & \textbf{Base ($\pi_0$)} & \textbf{Instruct ($\pi$)} & \textbf{Difference ($\Delta$)} \\
\midrule
Recovery States ($s_R$) & 0.1500 & 0.5800 & $+0.4300$ \\
Control States ($s_C$) & 0.3800 & 0.9200 & $+0.5400$ \\
\midrule
\textbf{Matched Contrast ($D_{\text{recovery}}$)} & -- & -- & \textbf{$-0.1100$} \\
\bottomrule
\end{tabular}
\end{table}

Under the evaluated state-matched protocol, we did not observe evidence of a recovery-specific advantage for the Instruct checkpoint over the Base checkpoint. The estimated matched recovery-specific checkpoint-interface contrast was $-0.1100$, with a 95\% descriptive problem-level bootstrap interval of $[-0.240, +0.030]$ (10,000 resamples).

The Instruct checkpoint showed higher continuation success than Base in both recovery ($+0.4300$) and control ($+0.5400$) states, but the improvement was larger for controls. Consequently, aggregate benchmark gains do not automatically translate into a recovery-specific advantage.

\section{Sensitivity and Reproducibility Analysis}
We assess matching quality across the 20 matched pairs. The mean normalized weighted-L1 distance is $d_{\text{mean}} = 0.0360$, median is $d_{\text{median}} = 0.0360$, and maximum distance is $d_{\max} = 0.0360$. All 20 pairs (20/20) satisfy both the standard ($d \le 0.25$) and tight ($d \le 0.10$) thresholds \cite{austin2011introduction}.

Table \ref{tab:sensitivity} reports covariate balance metrics. Standardized Mean Differences (SMDs) for trajectory depth and remaining length are $+0.0000$, while token length SMD is $+0.1333$.

\begin{table}[htbp]
\caption{Matching Distance, Covariate Balance \& Threshold Sensitivity}
\label{tab:sensitivity}
\centering
\begin{tabular}{lccc}
\toprule
\textbf{Metric / Covariate} & \textbf{Raw Difference} & \textbf{SMD / Distance} & \textbf{Status} \\
\midrule
Trajectory Depth & $+0.0000$ & $+0.0000$ & Exact Match \\
Remaining Length & $+0.0000$ & $+0.0000$ & Exact Match \\
Token Length & $+2.0000$ & $+0.1333$ & Balanced \\
\midrule
Mean Pair Distance ($d_{\text{mean}}$) & -- & $0.0360$ & $20/20 \le 0.25$ \\
Max Pair Distance ($d_{\max}$) & -- & $0.0360$ & $20/20 \le 0.10$ \\
\bottomrule
\end{tabular}
\end{table}

Independent re-computation from sealed raw evidence \texttt{RAW\_NEURAL\_ROLLOUTS.jsonl} (SHA-256: \texttt{51b5a157d9e4...}) yields exact agreement ($0.0$ discrepancy across all metrics).

\section{Limitations}
This study has explicit methodological scope boundaries:
\begin{itemize}
    \item \textbf{Model Scope}: Evaluated on a single model family (\texttt{Qwen2.5-Math-1.5B}) across two released checkpoints.
    \item \textbf{Benchmark Scope}: Evaluation uses 20 GSM8K test items ($N=20$). Benchmark pretraining contamination cannot be excluded.
    \item \textbf{Sample Size}: 5 stochastic rollout continuations per state/policy represent evaluation samples, not independent model training replications.
    \item \textbf{No Causal Claim}: Results represent a descriptive contrast under state matching, not a causal claim regarding post-training mechanisms.
\end{itemize}

\section{Conclusion}
We presented \texttt{recovery\_eval}, a prospective state-matched evaluation framework for language-model reasoning diagnostics. By pairing verifier-defined recovery states with structural reference controls and maintaining primitive rollout governance, the framework enables fine-grained evaluation of reasoning policies without confounding baseline fluency gains.

\bibliographystyle{IEEEtran}
\bibliography{references}

\end{document}
